\noindent\textbf{课程与资料范围.}
本卷依据 Scott A. Hughes 教授的 MIT OpenCourseWare 8.962 \emph{General Relativity}, Spring 2020,
按公开 Lecture Guide 与课程进度表的第 1--25 讲主题独立重编.
原课程为研究生课程, 每周两次课, 每次90分钟.
第24、25讲未录视频, 公开Guide有后牛顿展开、黑洞微扰和数值相对论概要;
本卷补推仅限公开主题, 不称完整课堂实录.

\noindent\textbf{先修与阅读路径.}
读者应掌握多元微积分、线性代数、常微分方程和经典力学中的拉格朗日量、作用量原理;
电磁学和数值积分有助于完成作业; MIT原先修为18.03、18.06和8.07.
本卷不要求预先学过微分几何: 从洛伦兹变换、四矢量及对偶空间出发, 逐步建立张量、流形、度规、联络与曲率,
再把这些对象用于爱因斯坦方程、辐射、宇宙学、恒星和黑洞.

\begin{center}
\small
\begin{tabular}{@{}p{.07\textwidth}p{.10\textwidth}p{.48\textwidth}p{.25\textwidth}@{}}
\toprule
本卷章 & 原讲次 & 内容进度 & 原课程作业节点\\
\midrule
1 & 1--5 & 狭义相对论、张量、体积元与守恒律 & 第4讲: PS1\\
2 & 5--8 & 度规、等效原理、联络、输运与对称 & 第7讲: PS2\\
3 & 9--11 & 测地线、曲率、潮汐与 Bianchi 恒等式 & 第9讲: PS3\\
4 & 12--13 & 爱因斯坦方程及作用量推导 & 第12讲: PS4\\
5 & 14--17 & 线性引力、动态源与引力辐射 & 第14/16讲: PS5/6\\
6 & 18--19 & 宇宙学、红移与距离 & 第18讲: PS7\\
7 & 20--21 & 球对称恒星、紧致性与引力检验 & 第20讲: PS8\\
8 & 22--23 & 黑洞、视界与轨道 & 第22讲: PS9\\
9 & 24--25 & 后牛顿展开、黑洞微扰、空间时间分解 & 第24讲: PS10/11\\
\bottomrule
\end{tabular}
\end{center}
第5讲的应力能量内容在第1章, 联络内容接入第2章.
第10--12章分别收入PS1--3、PS4--7、PS8--11的题面和解答.

\noindent\textbf{作业、考试与缺口.}
2020 年版作业没有公开; OCW 作业页实际提供的是 Spring 2006 的 11 套作业和 PS1、PS3、PS4、PS7 的四份官方提示.
本卷保留原题编号, 收入 51 道可获得或由官方提示恢复要求的大题, 全部给出本稿独立解答;
其中 PS7 第3题由官方提示中的完整拉氏量和变分要求恢复, 不能称为取得了教材原页.
PS3 第5、6题和 PS7 第4、5题只有 Carroll 教材题号, 原题面未取得, 因而明确保留为四项缺口, 不编造题目或答案.
本课程大纲明确写明无考试, 不存在应收入而遗漏的本课程期中或期末卷.

\noindent\textbf{署名与修改.}
原资料署名Scott A. Hughes、MIT OpenCourseWare.
本卷为AI辅助中文独立重编, 补充推导、题解与TikZ图; MIT未审核或认可.
OCW资料及中文改编按CC BY-NC-SA 4.0保留署名、非商业与相同许可要求.
完整来源、第三方例外与修改说明见第13章及源码的\texttt{LICENSE-NOTICE.md}.
